% Typeset transcription of Marco.pdf.
% Compile with: pdflatex Marco.tex
% The diagrams in the handwritten source are represented schematically.

\documentclass[11pt,a4paper]{article}

\usepackage[T1]{fontenc}
\usepackage[utf8]{inputenc}
\usepackage{amsmath}
\usepackage{amssymb}
\usepackage{xcolor}
\usepackage[margin=1in]{geometry}
\usepackage[colorlinks=true,linkcolor=blue!50!black,urlcolor=blue!50!black]{hyperref}

\setlength{\parindent}{0pt}
\setlength{\parskip}{0.65em}

\definecolor{accent}{HTML}{7A2C83}
\newcommand{\ID}{\operatorname{ID}}
\newcommand{\ED}{\operatorname{ED}}
\newcommand{\highlight}[1]{\textcolor{accent}{\textbf{#1}}}

\title{Looking for a Semantic Core in LLM Representations}
\author{Lecture by Marco Baroni\\Notes by Lucia Domenichelli}
\date{2026 School on Analytical Connectionism\\Gothenburg, 26 August 2026}

\begin{document}

\maketitle

\begin{abstract}
These notes ask whether different neural networks share a semantic core and how
that core changes across layers, architectures, languages, and modalities. The
discussion combines intrinsic-dimensionality analysis, information imbalance,
cross-lingual and visual comparisons, and syntactic/semantic ablations.
\end{abstract}

\section{Motivation}

As humans, two very different sentences can convey the same meaning. Meaning is
core to human language, but it is latent:

\[
  \textbf{meaning}
  \quad\Longrightarrow\quad
  \substack{\text{core to human language}\\[2pt]\text{but latent}.}
\]

The central question is whether different neural networks capture the same
meaning structure. This motivates the \emph{Platonic Representation
Hypothesis} and two questions:

\begin{enumerate}
  \item Is there really a shared semantic core?
  \item What are its properties across networks, network layers, modalities,
        and inputs?
\end{enumerate}

\section{Emergence of a High-Dimensional Representation Space}

An input is mapped to hidden-state vectors, which form a geometric
representation space:

\[
  \text{input}
  \longrightarrow
  \boxed{\text{hidden-state vectors}}
  \longrightarrow
  \text{representation space}
  \longrightarrow
  \mathcal{M}.
\]

Here, $\mathcal{M}$ is treated as a data manifold. Its dimensionality can be
studied at two levels:

\begin{itemize}
  \item the \textbf{extrinsic dimension} $\ED$, determined by the ambient
        representation space;
  \item the \textbf{intrinsic dimension} $\ID$, the effective number of
        parameters needed to characterize the data manifold.
\end{itemize}

The intrinsic dimension evolves across layers:

\[
  \ID\!\left(h^{(1)}\right)
  \longrightarrow
  \ID\!\left(h^{(2)}\right)
  \longrightarrow \cdots \longrightarrow
  \ID\!\left(h^{(L)}\right).
\]

These manifolds are only locally linear and are generally nonlinear, so the
analysis should not assume global linearity. The notes mention the
\textsc{GRIDE} estimator (Denti et al., 2022), implemented in
\textsc{DADApy}, for estimating intrinsic dimension.

\subsection{Experimental setup}

The language-model study uses a random sample of sentences:

\begin{itemize}
  \item $10{,}000$ fixed-length token sequences from the Pile;
  \item a shuffled version of the corpus as a comparison.
\end{itemize}

For a causal language model, the hidden state of the final token is used as the
representation of the sentence. In the example, the extrinsic dimension is
$4096$, while the measured intrinsic dimension is much smaller:

\[
  \ID \ll \ED = 4096.
\]

Across layers, the handwritten plot shows the intrinsic dimension rising to a
pronounced peak, falling, and then increasing slightly near the end. This peak
emerges during training and is visible across multiple Pythia checkpoints.

\subsection{What happens at the intrinsic-dimension peak?}

The next question is whether the $\ID$ peak coincides with the emergence of
linguistic information. The probing tasks are grouped as follows:

\begin{itemize}
  \item \textbf{Surface properties}
    \begin{itemize}
      \item sentence length;
      \item word content.
    \end{itemize}
  \item \textbf{Deeper linguistic properties}
    \begin{itemize}
      \item bigram shift;
      \item coordination inversion;
      \item odd-man-out detection.
    \end{itemize}
\end{itemize}

The notes report no correlation between the $\ID$ peak and the surface
properties, but a correlation with the deeper probes.

\section{A Systematic Study of Representational Geometry}

A systematic study is useful because extracting representations manually with
model hooks becomes too complex. A gzip-based linear/compression perspective is
also mentioned as a point of comparison.

For a quantitative analysis of semantic information in deep representations,
the notes compare visual and textual data. Such representations can differ in
extrinsic dimension and random-seed initialization, so the focus is placed on
their \emph{similarity structure}. Three comparison methods are listed:

\begin{itemize}
  \item neighborhood overlap;
  \item information imbalance;
  \item representational similarity analysis.
\end{itemize}

\subsection{Information imbalance}

Information imbalance is a local, asymmetric measure. Let $r^A_{ij}$ be the
rank of point $j$ among the neighbors of point $i$ in representation space $A$,
and define $r^B_{ij}$ analogously for space $B$. The formula recorded in the
notes is

\[
  \Delta(A \to B)
  =
  \frac{2}{N}
  \sum_{\substack{i,j\\ r^A_{ij}=1}}
  \frac{r^B_{ij}}{N}.
\]

Equivalently, this asks where the nearest neighbor in $A$ is ranked in $B$.
The direction matters:

\[
  \Delta(A \to B) \neq \Delta(B \to A)
  \quad\text{in general}.
\]

The measure runs approximately from $0$ to $1$. A value near $0$ indicates
that the spaces preserve the same local neighborhoods; larger values indicate
less shared predictive structure. The notes also mention a data-processing
property (Del Tatto et al., 2024).

\section{Semantic Convergence Across Languages and Models}

\subsection{Translations}

Information imbalance is measured across relative network depth for translated
sentences. The handwritten comparison uses

\[
  \text{English} \longrightarrow \text{Italian}
  \qquad\text{and}\qquad
  \text{English} \longrightarrow \text{Hungarian},
\]

with an open-weight DeepSeek model. Both sketches are U-shaped over relative
depth, suggesting stronger convergence in intermediate layers. The notes also
flag unequal training exposure: the model has seen less Hungarian. Because the
curve labels and this side note are not fully consistent in the handwritten
sketch, the precise ordering of the two curves should not be inferred from the
drawing alone.

\subsection{Language prevalence}

For Qwen3-8B, the notes plot minimum information imbalance against the online
presence of several languages (Hungarian, Italian, Spanish, and German). The
recorded correlation is

\[
  \rho\!\left(\text{online presence},\,\Delta_{\min}\right) = -0.82.
\]

Thus, greater online presence is associated with lower minimum information
imbalance in this comparison, suggesting stronger cross-lingual convergence.
Relative depth makes it possible to compare the same pattern across models with
different numbers of layers.

\section{Visual and Multimodal Representations}

The same analysis can be applied to image representations. The notes contrast
two visual models:

\begin{itemize}
  \item \textbf{ImageGPT}, a causal model, whose curve has a minimum in the
        intermediate layers;
  \item \textbf{DINOv2}, whose representations become increasingly informative
        toward later layers and are intended to serve as features for many
        downstream tasks.
\end{itemize}

Both tied- and untied-weight settings are considered, and the model still
appears to learn the correspondence.

The notes then compare images with their descriptions in otherwise unrelated
visual and language models. One directional comparison is

\[
  \text{CLIP-Text} \longrightarrow \text{CLIP-Image}.
\]

The qualitative result is that language predicts vision better than vision
predicts language.

\section{Differential Syntactic and Semantic Encoding}

The final study separates syntax and semantics in LLM representations. It
constructs syntactic and semantic twins, summarizes them with corresponding
centroids, and performs ablations:

\[
\begin{array}{ccc}
  \text{syntax twins} && \text{semantic twins} \\
  \downarrow && \downarrow \\
  \text{syntax centroids} && \text{semantic centroids}
\end{array}
\]

\begin{itemize}
  \item semantic ablation;
  \item syntactic ablation;
  \item evaluation with symmetrized information imbalance.
\end{itemize}

A natural symmetrized comparison is

\[
  \Delta_{\mathrm{sym}}(A,B)
  = \frac{1}{2}
    \bigl(\Delta(A \to B) + \Delta(B \to A)\bigr).
\]

Probing then branches into syntactic and semantic retrieval. The notes record
that removing semantics has a larger effect than removing syntax, while also
observing that syntactic ablation disrupts semantic information. Syntax,
semantics, and the residual representation do not by themselves capture all of
the structure:

\[
  \left.
  \begin{array}{l}
    \text{syntax} \\
    \text{semantics} \\
    \text{residual}
  \end{array}
  \right\}
  \longrightarrow
  \text{not the whole representation}
  \longrightarrow
  \text{semantic core}.
\]

\section{Main Takeaways and Open Direction}

\begin{itemize}
  \item A semantic core appears in the central layers of the models.
  \item The intrinsic-dimension peak is related to deeper linguistic probes,
        not to the surface probes considered here.
  \item Information imbalance provides a local, directional way to compare
        representations across layers, models, languages, and modalities.
  \item Unequal data exposure can affect apparent cross-lingual convergence.
  \item Syntax and semantics are not fully independent: ablating syntax can
        disrupt semantic information.
\end{itemize}

The next goal is to map representations onto a shared semantic manifold that
supports composed meaning:

\[
  \text{model representations}
  \longrightarrow
  \text{\highlight{shared semantic manifold}}
  \longrightarrow
  \text{composed meaning}.
\]

The closing question is how simple concepts can be retrieved from this
manifold. The proposed direction is \highlight{polar probes}.

\section*{Transcription Note}

This document is a cleaned transcription of the handwritten source. Obvious
spelling and grammatical slips have been normalized. Hand-drawn plots are
described qualitatively because the source provides no numerical coordinates;
page 9 of the PDF is blank.

\end{document}
